\BeleskeID{09_2-s019}
% element: 09_2-s019:e04
Porastom ulaznog napona izlazni napon opada, dok napon između sorsa i drejna P tranzistora raste. Uslovi za istovremeno zasićenje u oznakama priključaka su
\[V_{DS,N}=V_O\ge V_{GS,N}-V_{Tn}=V_I-V_{Tn}\]
\[V_{DS,P}=V_O-V_{DD}\le V_{GS,P}-V_{Tp}=-V_{DD}-V_{Tp}\]
Tada i tranzistor T1 i tranzistor T2 imaju uslove za zasićenje. Polazna ravnoteža struja je
\[\frac{k_p}2(V_{GS,P}-V_{Tp})^2=\frac{k_n}2(V_{GS,N}-V_{Tn})^2\]
U idealizovanom modelu dobija se ulazni napon $V_\infty$, pri kojem izlaz više nije jednoznačno određen zanemarenim efektima. Model sa kratkim kanalom i modulacijom dužine kanala vraća zavisnost od izlaznog napona.
